% Abstract text for CsInH3 paper. Include via \input{} inside \begin{abstract}...\end{abstract} in main.tex.

Hydrogen-rich superconductors achieve critical temperatures above 200~K,
but only at pressures exceeding 150~GPa --- conditions confined to
diamond anvil cells and incompatible with practical applications or bulk
characterization. Here we screen ternary perovskite hydrides MXH$_3$
(M~=~Cs, Rb, K; X~=~In, Ga) using density functional theory, density
functional perturbation theory, isotropic Eliashberg theory, and
stochastic self-consistent harmonic approximation (SSCHA) anharmonic
corrections. The top-ranked candidate, CsInH$_3$ in the cubic
Pm$\bar{3}$m structure, achieves $T_c = 214$~K at $P = 3$~GPa
($\mu^* = 0.13$, $\lambda_{\mathrm{anh}} = 2.26$) after SSCHA
renormalization --- matching the experimental $T_c$ of H$_3$S at a
pressure $30\times$ lower. CsInH$_3$ is thermodynamically stable
($E_{\mathrm{hull}} = 6$~meV/atom) and quantum-stabilized at 3~GPa by
hydrogen zero-point motion. Despite this result, 300~K
superconductivity is not achievable in the MXH$_3$ family due to the
intrinsic $\lambda$--$\omega_{\mathrm{log}}$ tradeoff that caps $T_c$
near $\sim 230$~K. These findings identify perovskite hydrides as a
route to high-$T_c$ superconductivity at pressures accessible to
large-volume press synthesis.
